\subsection{章末確認問題}
\begin{enumerate}[leftmargin=2em]
\item 工学プロセスの六つの段階を、自分の言葉で説明せよ。
\item 工学設計が一つの正解を探す活動ではない理由を説明せよ。
\item 抽象化とカプセル化の違いを、APIを例に説明せよ。
\item 計画実験、観察研究、回顧研究はどのように使い分けるか。
\item 分析単位、母集団、標本の違いを、利用者満足度調査を例に説明せよ。
\item 相関が因果を意味しない理由を述べよ。
\item 名義尺度、順序尺度、間隔尺度、比率尺度で許される演算の違いを説明せよ。
\item CMMI成熟度段階を整数で表したときに起こり得る測定理論上の問題を説明せよ。
\item GQMが「測りやすいものを測る」発想と異なる理由を説明せよ。
\item 標準が設計制約になる場合の例を挙げよ。
\item 根本原因分析と表面的な対症療法の違いを説明せよ。
\item インダストリー4.0でソフトウェア工学が重要になる理由を説明せよ。
\end{enumerate}
\begin{notebox}{解答の方向性}1は真の問題、基準、候補解、評価、選択、監視を含める。2は制約下で複数解を比較する点を述べる。3は細部を隠して水準を定める抽象化と、インタフェース越しに内部を隠すカプセル化を区別する。7と8では、数値であっても尺度が違えば平均や乗除が不適切になる点を説明する。10から12では、標準、再発防止、接続された製造システムという語を使うとよい。\end{notebox}
